\begin{textsample}{POS Dim 1 – vem\_ed\_25 – Score 17.00 – vem\_ed\_25/t134.txt}  \label{ex:f1_pos_071}
4º Congresso Red LatAM COIL: diversidade, equidade e inclusão Workshop apresentado por Osvaldo Succi Junior (CPS / Brasil), José Luis Jiménez (Universidad Católica Andrés Bello, Venezuela) e Mirjam Hauck (The Open University, Reino Unido) Entre 4 e 6 de setembro de 2024, ocorreu o 4º Congresso da Red LatAM COIL, com o tema “diversidade, equidade e inclusão”.
Um dos objetivos da rede \textbf{é} \textbf{expandir} os benefícios da abordagem COIL (Collaborative Online International Learning) como \textbf{estratégia} de internacionalização do currículo na educação \textbf{superior}, envolvendo países da América Latina e outras \textbf{partes} do mundo.
Sessões paralelas, painéis de discussão e workshops compuseram o evento on-line, \textbf{que} \textbf{contou} com as convidadas especiais Hope Windle (SUNY COIL Center) e Rosi León (DePaul University).
Na abertura, após a palavra da presidenta da Red LatAM COIL, Brenda García Portillo, apresentaram-se os membros do conselho executivo - entre eles, Ana Carolina Freschi, da equipe dos PCIs / Cesu.
O tema da conferência de abertura proferida por Hope Windle, diretora do SUNY COIL Center, \textbf{foi} “COIL and diversity, equity and inclusion: beyond icebergs and stereotypes to community building”.
Hope Windle ressaltou \textbf{que} diversidade \textbf{começa} pelo esforço dos professores em pronunciar corretamente o nome dos alunos de várias origens.
Celebrou também a translinguagem (spanglish, portuñol) e \textbf{defendeu} sessões de quebra-gelo (icebreakers) em \textbf{que} se desconstroem estereótipos.
Apresentou o acrônimo TTLC (time, trust, leadership, communication) como caminho para a inclusão, despertando empatia e consideração pela diversidade de pensamento.
No segundo dia da conferência, Osvaldo Succi Junior, coordenador dos PCIs / Cesu, apresentou, com José Luis Jiménez (Universidad Católica Andrés Bello, Venezuela) e Mirjam Hauck (The Open University, Reino Unido) o workshop “Humanizing STEM through COIL and Critical Virtual Exchange: Empowering Educators to Foster Interdisciplinary and Culturally Inclusive Learning”. \textbf{continuação} A \textbf{proposta} de humanização do STEM (Science, Technology, Engineering, Mathematics) por meio dos projetos COIL passa pelo alinhamento com os Objetivos de Desenvolvimento \textbf{Sustentável} (ODS) da ONU.
Três \textbf{estratégias} \textbf{são} \textbf{possíveis}: colaboração entre cursos STEM-STEM – exemplos: química, ajudando comunidades sul-africanas a \textbf{produzir} sabões artesanais para melhorar a renda familiar; matemática, \textbf{analisando} dados da migração venezuelana para compreender os desafios dos migrantes; desenvolver projetos \textbf{que} envolvam disciplinas de humanidades e ciências “duras” para abordar desafios sociais – exemplo: ciência da computação e jornalismo, desenvolvendo jogos \textbf{que} despertam consciência sobre \textbf{paz}, justiça e direitos \textbf{humanos}; envolver ONGs dos dois países participantes do projeto \textbf{colaborativo} e \textbf{propor} soluções para as comunidades – exemplo: projeto entre universidades do Brasil e da Venezuela \textbf{que} colaboraram com ONGs para lidar com a remoção forçada de comunidades indígenas devido à extração mineral e ao desmatamento na Bacia Amazônica.
Osvaldo Succi Junior sugere, nesse \textbf{contexto} de humanização dos projetos, uma nova \textbf{proposta} de PCI / COIL.
Em vez de desenhar o projeto com base em \textbf{disciplinas}, \textbf{partir} de um ou mais ODS para a solução de problemas \textbf{locais} ou globais.
Na \textbf{parte} \textbf{prática} do workshop, os participantes esboçaram um projeto COIL contemplando o ODS 3 (Saúde e Bem-Estar) e o 10 (Redução das Desigualdades).
Tarefa principal: desenvolver campanhas de comunicação para desconstruir mitos relacionados às vacinas no México e nas Bahamas, em colaboração com ONGs como Médicos sem Fronteiras e Cruz Vermelha.

% matched lemmas: analisar, colaborativo, começar, contar, contexto, continuação, defender, disciplina, estratégia, expandir, humano, local, parte, partir, paz, possível, produzir, propor, proposta, prático, que, ser, superior, sustentável
\end{textsample}
